\section{Introduction}
\label{sec:intro}

Choosing the fastest collective-communication strategy for a
distributed-training workload is challenging, and
two prevailing methods have substantial limitations. Running each
candidate strategy end-to-end on the actual workload is direct
but costly: realistic GPU-cluster experiments are expensive,
and the cost grows with the number of strategies tried.
Ranking strategies by a cheap microbenchmark on small
isolated calls is much cheaper, but on a Trainium cluster we
show that multiple strategies that win microbenchmarks are
materially \emph{slower} at actual training runs than a
different strategy that loses in microbenchmarks
(Figure~\ref{fig:disagree}). We call this \emph{cross-scope
inversion}: moving an operation from an isolated microbenchmark
into a real multi-rank training step changes the per-call cost,
because the surrounding training graph changes how the runtime
launches, fuses, and lays out tensors for the call.

A third approach, scoring strategies against a cost-model
simulator, has been explored in prior work
(SimAI~\cite{wang2025simai},
ASTRA-sim~\cite{rashidi2020astrasim,won2023astrasim2}). But
existing simulators are written and tuned against a fixed
reference configuration: the numbers they use internally ---
per-collective bandwidth, the overhead each time the runtime
launches a graph, when adjacent calls fuse, etc. --- drift
across machines, model shapes, and library versions,
and re-tuning them by hand for each new environment
or training task is again costly. The
simulator-based approach gets even harder on the closed-source
vendor stacks production training increasingly runs on (AWS
Trainium, Google TPU, Meta MTIA), because of two specific
properties of these stacks.
\textbf{(C1) No visibility into the runtime:} when an operation
(e.g., \texttt{xm.all\_to\_all} from Trainium) is unexpectedly slow, developers
cannot inspect the vendor library to see why --- its internals are
not source-available, and there are limited per-primitive 
configuration parameters to calibrate a simulator against.
\textbf{(C2) No low-level control:} the vendor's collective
primitives (\texttt{xm.all\_gather}, \texttt{xm.all\_to\_all},
\dots) are called through a fixed API; there is no exposed
schedule-level intermediate representation
--- which, in the open-stack systems
MSCCL~\cite{cowan2022msccl} and
TACCL~\cite{shah2023taccl} target, is the surface that lets a
user choose how chunks are split, routed through the network,
or fused inside a collective.

This paper introduces a new design that addresses each of
these challenges.
\textbf{For (C1)}, we approximate visibility into the closed
runtime by having an LLM agent build a cost-model simulator
on the target deployment itself: the LLM writes probing code
that measures the simulator's internal numbers against the
actual device, model shapes, and library version the workload
will use. This does not require internal visibility into the
runtime, but it gives a usable proxy for ``is this strategy
fast for this specific task?'' whose estimates re-adapt to the
deployment when any of the device, model, or library changes
--- so the simulator stays workload-faithful where a
hand-coded one would drift.
\textbf{For (C2)}, we do not try to recover low-level control.
We accept the vendor API as fixed and search one layer above
it --- over which primitives to call, in what order, and how
to reshape, pad, or slice the input/output tensors around each
call. The LLM proposes strategies at this layer; the
workload-calibrated simulator scores them; top-scoring
strategies are validated on real hardware before deployment.

One might ask: if the simulator's own internal numbers come
from on-device probes, isn't this just the microbenchmark
option in a different wrapper? Our ablation
(\S\ref{sec:ablation-nosim}) is set up to answer exactly that.
The difference is in \emph{what} we probe and \emph{how} we
probe it. Our simulator encodes hardware-specific performance
models, and the LLM is guided to calibrate them using probe
workloads that resemble real model-training settings (real
tensor shapes, full step boundaries, multi-rank dispatch
patterns), rather than the small isolated calls a standalone
microbenchmark would measure. The probes therefore capture
behavior that small microbenchmarks miss, while remaining far
cheaper than running a full training harness end-to-end on
each candidate strategy. When we ablate the simulator and let
the LLM rank strategies by its own small microbenchmarks
instead, the $1.40\times$ OLMoE-10B speedup collapses to
$1.00\times$: the LLM converges to structurally the same
strategy practitioners already deploy, because microbench
rankings are exactly what selected that strategy in the first
place. In other words, the work-load-faithful simulator, not
the LLM, is what picks the strategy that wins at actual
training runs.


\paragraph{Closed-stack collective communication library} A growing class of
production training accelerators operates on proprietary collective communication libraries. AWS Trainium~\cite{aws_trainium_neuron,
aws_trainium2_2024} --- 30--40\% better
price-performance than NVIDIA H100 on transformer workloads and
adopted at scale for Mixture-of-Expert (MoE) and Llama training --- ships a Neuron
Runtime that exposes collectives only through the PyTorch/XLA
integration (XLA = Accelerated Linear Algebra, PyTorch's
compiler bridge to non-NVIDIA accelerators) --- i.e.\ as
\texttt{xm.*} entry points such as
\texttt{xm.\allowbreak{}all\_gather} --- behind which the Neuron
compiler emits NEFF binaries (Neuron Executable File Format;
the compiled artifact loaded onto each NeuronCore) from
unpublished sources. There is no
user-configurable schedule layer: per-collective chunking and
network-interface-card (NIC) routing cannot be reached from
above the runtime. Google TPU and Meta MTIA exhibit the same
constraint: the programmable schedule surface that MSCCL- and
TACCL-style synthesis depend on is not exposed. The schedule-synthesis paradigm
of the open-stack era cannot be ported here in any straightforward
way. Above the \texttt{xm.*} (or equivalent vendor) boundary, the
design choice that remains is which primitive to call, in what
order, and against what tensor layout. We refer to this as the \emph{strategy} layer,
in contrast to the \emph{schedule} layer that MSCCL/TACCL
operate on, which decides how each collective message is split
into chunks, which network paths each chunk takes, and how the
chunks are pipelined across devices. The strategy layer is less expressive
but, as we show, contains strategies $1.4$--$3.2\times$ faster
than the developer baseline.

\paragraph{Related schedule-synthesis work and why it does not
apply here.} Automated synthesis of collective-communication
algorithms has been studied at the \emph{schedule} layer ---
which message chunks move along which paths, in what order, against
which network adapters. MSCCL and MSCCLang~\cite{cowan2022msccl}
emit XML or DSL schedules for the NCCL/RCCL runtime~\cite{nvidia_nccl} to consume;
TACCL~\cite{shah2023taccl} uses a sketch language and an SMT
solver; AutoCCL~\cite{xu2025autoccl} adds tuning;
SCCL~\cite{cai2021sccl} synthesizes Pareto-optimal algorithms
from first principles. All share an infrastructural prerequisite:
the underlying runtime exposes a programmable intermediate
representation (IR) for the schedule.


\paragraph{Validation and contributions.}
We validate our design on the eight collective problems listed
in Table~\ref{tab:problems} --- four that arise in
Mixture-of-Expert (MoE) training~\cite{shazeer2017outrageously}
and four that arise in dense-transformer (Llama-style) training
--- on a 7-node \texttt{trn1.32xlarge} cluster (224 ranks).

\begin{table}[h]
\centering\scriptsize
\setlength{\tabcolsep}{3pt}
\setlength{\abovecaptionskip}{2pt}
\setlength{\belowcaptionskip}{-2pt}
\begin{tabular}{@{}p{0.28\columnwidth}p{0.67\columnwidth}@{}}
\toprule
\textbf{Problem} & \textbf{Role in the training step} \\
\midrule
\multicolumn{2}{@{}l}{\textit{MoE-side}} \\
AllToAllV~\cite{lepikhin2021gshard}      & MoE token dispatch/combine; counts vary \\
Uniform AllToAll~\cite{zhou2022expertchoice} & MoE dispatch/combine; counts equal \\
Ring KV~\cite{liu2024ringattention}      & rotate KV shards around the rank ring \\
Distributed CE~\cite{shoeybi2019megatron} & loss over rank-sharded vocabulary \\
\midrule
\multicolumn{2}{@{}l}{\textit{Llama-side}} \\
PP cross-stage~\cite{huang2019gpipe}     & activations/grads between PP stages \\
TP MLP~\cite{shoeybi2019megatron}        & all-reduce after MLP in tensor parallel \\
FSDP weight prefetch~\cite{zhao2023pytorchfsdp} & gather sharded weights per layer \\
Layer-block AR~\cite{korthikanti2023reducing} & all-reduce grads at layer blocks \\
\bottomrule
\end{tabular}
\caption{The eight collective problems we evaluate.}
\label{tab:problems}
\end{table} The baseline we measure against is the
\emph{internal-AWS-optimized strategy} AWS uses in production
today to train large-scale MoE and dense transformer models on
Trainium, built and maintained by five experienced AWS
developers. On OLMoE-10B the deployed agent strategy delivers
\textbf{1.40$\times$} steady-state step-time speedup with
matched descent on a 2500-step
real-OpenWebText~\cite{Gokaslan2019OpenWeb} AdamW run
(Figure~\ref{fig:loss-curve}); on Llama-7B it delivers
\textbf{3.24$\times$}. Per-primitive per-step speedups range
$4$--$12\times$, above the end-to-end ratio because constant
non-collective compute dilutes the headline.

The contributions are
(i)~to our knowledge the first system that finds faster
collective-communication strategies on top of a closed-source
vendor library, delivering \textbf{1.40$\times$} end-to-end on
OLMoE-10B and \textbf{3.24$\times$} on Llama-7B over the
production baseline;
(ii)~a cost-model simulator whose per-term constants are
re-fit per-environment by LLM-written on-device probes,
keeping the simulator workload-faithful without
hand-coding new probe scripts for each cluster; and
(iii)~the closed-stack search loop combining the above with
LLM-driven mutation and hardware validation, with code
emitted as drop-in runtime modules and algorithm-equivalent
descent verified on real tokens.

\paragraph{Scope and generalization.} We instantiate our design
on AWS Trainium. The six cost-model terms are Trainium-specific,
but the categories they fall into (per-graph-launch overhead,
per-collective bandwidth ceiling, device-vs.-host copy cost,
runtime cache-reload cost) are general to closed-source vendor
runtimes that compile their own collective binaries. Google TPU
and Meta MTIA each expose closed equivalents; we expect the
strategy-search design to port with the cost-model term set
re-calibrated, and leave that to follow-up work.
